% META: {"id": "TSPE-DERCONV-EX-033", "chapitre": "TSPE-DERIVATION-CONVEXITE", "type_objet": "exercice", "capacites": ["TSPE-DERCONV-C4"], "capacites_codes": ["C4"], "methodes": ["M4"], "parcours": 3, "competences": ["modeliser", "raisonner"], "duree_min": 20, "mode_creation": "ex_nihilo", "sources_inspiration": [], "fichier_tex": "chapitres/TSPE-DERIVATION-CONVEXITE/exercices/TSPE-DERCONV-EX-033.tex", "corrige_tex": "chapitres/TSPE-DERIVATION-CONVEXITE/corriges/TSPE-DERCONV-CO-033.tex", "status": "approved"}
% BEGIN-VERIFY
% from sympy import *
% x = symbols('x')
% # f''(x) = 12x^2 - 24, f'(0)=4, f(0)=1
% fpp = 12*x**2 - 24
% fp = integrate(fpp, x) + 4  # f'(0)=4
% assert expand(fp) == 4*x**3 - 24*x + 4
% f = integrate(fp, x) + 1  # f(0)=1
% assert f == x**4 - 12*x**2 + 4*x + 1
% assert solve(fpp, x) == [-sqrt(2), sqrt(2)]
% assert solve(fp, x)  # 3 racines reelles
% crit = solve(fp, x)
% assert len(crit) == 3
% END-VERIFY
\begin{exercice}{TSPE-DERCONV-EX-033}{3}{20}
On donne $f''(x) = 12x^2 - 24$, $f'(0) = 4$ et $f(0) = 1$.
\begin{enumerate}
  \item Determiner l'expression de $f$.
  \item Etudier la convexite et identifier les points d'inflexion.
  \item Esquisser la courbe en determinant les extremums.
\end{enumerate}

\end{exercice}
